% paper/sections/10b_implementation_details.tex
% §10 続き：ブートストラップ処理・DCCD パラメタ設計

% ─────────────────────────────────────────────────────────────────────────────
\subsection{ブートストラップ処理：初期化シーケンス}
\label{sec:bootstrap}
% ─────────────────────────────────────────────────────────────────────────────

シミュレーション開始前に，以下の初期化シーケンス（ブートストラップ処理）を一度だけ実行する．
格子生成には Level Set 場 $\phi^{(0)}$ が必要だが，$\phi^{(0)}$ の評価には格子が必要という
循環参照を，一様格子による仮評価で解消する．

\begin{tcolorbox}[algbox, title={ブートストラップ処理：初期化シーケンス}]
\begin{enumerate}
  \item \textbf{一様格子上での $\phi^{(0)}$ 計算：}
        $N_x \times N_y$ 等間隔格子上で $\phi^{(0)}_{i,j}$ を計算する
        （界面形状が解析的に与えられる場合は直接代入；そうでない場合は
        一様格子上の符号付き距離関数として初期化）．
  \item \textbf{非一様格子の生成（§~\ref{sec:grid_2d}）：}
        $\phi^{(0)}$ から界面適合非一様格子を生成する：
        \begin{enumerate}
          \item $x$ 方向の密度関数 $\omega^x_i = 1 + (\alpha-1)\,\varepsilon_g\,\delta_\varepsilon^*(\bar{\phi}^x_i)$ を計算し，
                §~\ref{sec:grid_gen} のステップ2--5で物理座標 $\{x_i\}$ とメトリクス係数 $\{J_x|_i\}$ を得る．
          \item $y$ 方向も同様に $\{y_j\}$ と $\{J_y|_j\}$ を得る．
        \end{enumerate}
  \item \textbf{CLS 変数の初期化：}
        非一様格子上で $\psi^{(0)}_{i,j} = H_\varepsilon(\phi^{(0)}_{i,j})$ を評価し，
        必要に応じて CLS 再初期化（§~\ref{sec:reinit}）を数ステップ実行して
        $|\bnabla\phi| \approx 1$ を確保する．
  \item \textbf{CCD 演算子の構築：}
        非一様格子のメトリクス係数を用いて CCD ブロック三重対角行列
        （§~\ref{sec:ccd_matrix}）を組み立て，LU 分解を事前計算する．
  \item \textbf{初期圧力場の設定：}
        静的平衡条件 $p^{(0)} = \sigma\kappa_0$（Young--Laplace 圧力）を初期値として設定する．
        曲率 $\kappa_0$ は $\psi^{(0)}$ から CCD 微分を用いて評価する．
\end{enumerate}

\textbf{出力：}
非一様格子 $\{(x_i, y_j)\}$，メトリクス係数 $\{J_x|_i, J_y|_j\}$，
CLS 場 $\psi^{(0)}$，CCD 演算子（LU 分解済み），初期圧力場 $p^{(0)}$．
\end{tcolorbox}

格子の時間更新モード（Mode 1: 固定格子，Mode 2: 再生成）については
§~\ref{sec:grid_ale} を参照されたい．

% ─────────────────────────────────────────────────────────────────────────────
\subsection{DCCD パラメタの設計と適応制御}
\label{sec:dccd_params}
% ─────────────────────────────────────────────────────────────────────────────

§~\ref{sec:dccd_filter_theory} で導入した DCCD フィルタの理論的枠組みに基づき，
本節では散逸強度 $\varepsilon_{d,\max}$ の具体的な設計値と，
CLS 変数 $\psi$ に基づく適応制御アルゴリズムを与える．

% ─────────────────────────────────────────────────────────────────────────────
\subsubsection{散逸強度 \texorpdfstring{$\varepsilon_{d,\max}$}{ε\_d,max} のスペクトル設計}
\label{sec:dccd_eps_design}
% ─────────────────────────────────────────────────────────────────────────────

スイッチ関数が決まれば，次は散逸強度の上限 $\varepsilon_{d,\max}$ を設計する問題となる．
設計の指針は以下の 2 つの要求の\textbf{同時達成}である：
\begin{enumerate}[label=(\alph*)]
  \item \textbf{精度保存：} 低波数域（長波長の物理情報）は CCD の $\Ord{h^6}$ 精度を損なわない．
  \item \textbf{高周波抑制：} ナイキスト近傍（$\xi \gtrsim 2\pi/3$，波長 $\lesssim 3h$）を
        有意に減衰させる．
\end{enumerate}

\medskip
\noindent\textbf{定量的設計．}

フィルタ伝達関数 $H(\xi;\varepsilon_d) = 1 - 4\varepsilon_d\sin^2(\xi/2)$ において，
低波数（$\xi \leq \xi_L$）での精度要求を：
\begin{equation}
  |H(\xi_L;\varepsilon_d) - 1| = 4\varepsilon_d\sin^2\!\left(\frac{\xi_L}{2}\right) \leq \delta_\mathrm{acc}
  \label{eq:accuracy_req}
\end{equation}
高波数（$\xi = \pi$）での減衰要求を：
\begin{equation}
  H(\pi;\varepsilon_d) = 1 - 4\varepsilon_d \leq 1 - D_\mathrm{target}
  \quad\Leftrightarrow\quad
  \varepsilon_d \geq \frac{D_\mathrm{target}}{4}
  \label{eq:damping_req}
\end{equation}
と定式化する．典型的な二相流計算では，$\xi_L = \pi/3$（波長 $\geq 6h$ を精度保存），
$\delta_\mathrm{acc} = 0.05$（5\% 以内の振幅変化を許容），
$D_\mathrm{target} = 0.20$（ナイキスト成分を 20\% 減衰）とする：

\begin{tcolorbox}[resultbox, title={$\varepsilon_{d,\max}$ の設計値と設計根拠}]
\label{box:eps_design}
精度要求（式~\eqref{eq:accuracy_req}，$\xi_L = \pi/3$，$\delta_\mathrm{acc}=0.05$）から：
\begin{equation}
  4\varepsilon_d\sin^2(\pi/6) = 4\varepsilon_d\cdot\frac{1}{4} = \varepsilon_d \;\leq\; 0.05
  \label{eq:eps_bound_accuracy}
\end{equation}

減衰要求（式~\eqref{eq:damping_req}，$D_\mathrm{target}=0.20$）から：
\begin{equation}
  \varepsilon_d \;\geq\; \frac{0.20}{4} = 0.05
  \label{eq:eps_bound_damping}
\end{equation}

両制約が $\varepsilon_{d,\max} = 0.05$ で\textbf{等号達成}：

\begin{equation}
  \boxed{\varepsilon_{d,\max} = 0.05}
  \label{eq:eps_max}
\end{equation}

\textbf{解釈：}
波長 $6h$ 以上（物理的に意味ある成分）の精度劣化は最大 5\%，
ナイキスト成分（波長 $2h$）は 20\% 減衰する．
より強い安定化が必要な場合（高密度比・高 Reynolds 数計算）は
$\varepsilon_{d,\max} = 0.10$〜$0.20$ まで拡張できるが，
この場合は波長 $6h$ の精度劣化も 10〜20\% に増加することに注意する．

\medskip
\textbf{安定性制約の確認：}
$\varepsilon_{d,\max} = 0.05 \leq 0.25$（式~\eqref{eq:dccd_transfer} の安定限界 $1/4$）を満足する．$\square$
\end{tcolorbox}

\noindent\textbf{フィルタ伝達関数の波数プロファイル（$\varepsilon_d = 0.05$）：}
\begin{align*}
  H(0;\, 0.05)     &= 1.00 \quad (\xi = 0,\text{ DC 成分，変化なし}) \\
  H(\pi/3;\, 0.05) &= 1 - 4(0.05)\sin^2(\pi/6) = 1 - 0.05 = 0.95 \quad (\text{波長}\ 6h,\ 5\%\ \text{減衰}) \\
  H(2\pi/3;\, 0.05)&= 1 - 4(0.05)\sin^2(\pi/3) = 1 - 0.15 = 0.85 \quad (\text{波長}\ 3h,\ 15\%\ \text{減衰}) \\
  H(\pi;\, 0.05)   &= 1 - 4(0.05) = 0.80 \quad (\text{波長}\ 2h,\ 20\%\ \text{減衰})
\end{align*}

% ─────────────────────────────────────────────────────────────────────────────
\subsubsection{適応散逸制御則と実装アルゴリズム}
\label{sec:dccd_adaptive}
% ─────────────────────────────────────────────────────────────────────────────

\textbf{適応散逸強度．}
格子点 $i$ での散逸強度を CLS 変数 $\psi_i^n$（時刻 $t^n$）によって点別に制御する：
\begin{equation}
  \varepsilon_d^{(i)} \;=\; \varepsilon_{d,\max} \cdot S(\psi_i^n)
  \;=\; \varepsilon_{d,\max}\,(2\psi_i^n - 1)^2
  \label{eq:adaptive_eps}
\end{equation}
これにより：
\begin{itemize}
  \item $\psi_i \approx 0.5$（界面）：$\varepsilon_d^{(i)} \approx 0$，フィルタ無効，標準 CCD に退化
  \item $\psi_i \approx 0$ または $1$（バルク）：$\varepsilon_d^{(i)} \approx \varepsilon_{d,\max}$，
        最大散逸でノイズを積極的に排除
\end{itemize}

\begin{tcolorbox}[algbox, title={適応 Dissipative-CCD フィルタの実装アルゴリズム}]
\label{alg:dccd}
\textbf{入力：} スカラー場 $f$ の格子値 $\{f_i\}$，CLS 変数 $\{\psi_i^n\}$，散逸強度 $\varepsilon_{d,\max}$

\begin{enumerate}
  \item \textbf{CCD 求解：}
        標準 CCD ブロック Thomas アルゴリズム（§~\ref{sec:ccd_matrix}）を実行し，
        $\{f'_i\}$，$\{f''_i\}$ を得る．
  \item \textbf{点別散逸強度の計算：}
        各格子点 $i$ に対して
        \[
          \varepsilon_d^{(i)} = \varepsilon_{d,\max}\,(2\psi_i^n - 1)^2
        \]
        を計算する（計算コスト：$O(N)$，数値演算 1 回/点）．
  \item \textbf{1 次微分のフィルタ適用：}
        \[
          \tilde{f}'_i = f'_i + \varepsilon_d^{(i)}\,(f'_{i+1} - 2f'_i + f'_{i-1}),
          \quad i = 1, \ldots, N-2
        \]
        境界点 $i=0, N-1$：周期 BC では折り返しインデックスを使用；
        壁面・流出 BC では境界点をフィルタ適用前の CCD 値のまま保持する（§~\ref{sec:dccd_bc} 参照）．
  \item \textbf{2 次微分のフィルタ適用：}
        \[
          \tilde{f}''_i = f''_i + \varepsilon_d^{(i)}\,(f''_{i+1} - 2f''_i + f''_{i-1}),
          \quad i = 1, \ldots, N-2
        \]
  \item \textbf{後続計算への引き渡し：}
        フィルタ適用済みの $(\tilde{f}', \tilde{f}'')$ を曲率計算・速度勾配評価・
        圧力勾配計算（第~\ref{sec:pressure}章）に用いる．
\end{enumerate}

\textbf{2 次元への拡張：}
2 次元計算では $x$ 方向と $y$ 方向に独立に 1 次元フィルタを適用する（次元分割法）：
\[
  \tilde{f}'_{x,ij} = f'_{x,ij} + \varepsilon_d^{(ij)}\,(f'_{x,i+1,j} - 2f'_{x,ij} + f'_{x,i-1,j}),
  \quad
  \tilde{f}'_{y,ij} = f'_{y,ij} + \varepsilon_d^{(ij)}\,(f'_{y,i,j+1} - 2f'_{y,ij} + f'_{y,i,j-1})
\]
$f''_{xx}, f''_{yy}, f''_{xy}$ に対しても同様に適用する．

\textbf{計算コスト：} ステップ 1（CCD，$O(N^2)$ for 2D）+ ステップ 2--4（フィルタ，$O(N^2)$）．
フィルタは CCD の後処理として追加コストが小さい．
また，$\tilde{f}'$ と $\tilde{f}''$ はフィルタ段階で\textbf{独立}に補正する——
すなわち $\tilde{f}''$ は $\tilde{f}'$ から再連立せずに計算する（計算コスト削減のための設計選択）．
\end{tcolorbox}
